\clearpage\pagestyle{otMain}\fontsize{11}{15}\selectfont\corrigesChapitre{16}{Espace : sphère, boule et sections}
\begin{blocLivre}{cor-16-1}\begin{corrige}{1}{Un pavé}\demarche{Multiplier les trois dimensions, puis convertir.}
$V=8\times5\times3=\res{120\text{ cm}^3}$. Comme $1\text{ L}=1000\text{ cm}^3$, cela fait $\res{0{,}12\text{ L}}$.
\label{fig-cor-16-1}\figureGeo[corrige]{espace-01}
\end{corrige}\end{blocLivre}
\begin{blocLivre}{cor-16-2}\begin{corrige}{2}{Un cylindre}\demarche{Calculer l’aire du disque de base.}
$V=\pi\times3^2\times10=\res{90\pi\text{ cm}^3}\approx\res{283\text{ cm}^3}$. Le rayon est déjà fourni, il ne faut pas le diviser par deux.
\label{fig-cor-16-2}\figureGeo[corrige]{espace-02}
\end{corrige}\end{blocLivre}
\begin{blocLivre}{cor-16-3}\begin{corrige}{3}{Une pyramide}\demarche{Utiliser la hauteur perpendiculaire à la base.}
La base est un carré d’aire 36 cm$^2$. Le volume vaut $36\times9/3=\res{108\text{ cm}^3}$. La division par trois distingue la pyramide du prisme de même base et de même hauteur.
\label{fig-cor-16-3}\figureGeo[corrige]{espace-03}
\end{corrige}\end{blocLivre}
\begin{blocLivre}{cor-16-4}\begin{corrige}{4}{Un cône}\demarche{Passer du diamètre au rayon.}
Le rayon est 4 cm. $V=\pi\times4^2\times9/3=\res{48\pi\text{ cm}^3}$. Employer 8 comme rayon multiplierait le résultat par quatre.
\label{fig-cor-16-4}\figureGeo[corrige]{espace-04}
\end{corrige}\end{blocLivre}
\begin{blocLivre}{cor-16-5}\begin{corrige}{5}{Une boule}\demarche{Élever le rayon au cube, puis multiplier par $4\pi/3$.}
$V=\dfrac43\pi\times27=\res{36\pi\text{ cm}^3}\approx\res{113\text{ cm}^3}$. Un grand cercle a pour rayon 3 cm : sa longueur vaut $2\pi\times3=\res{6\pi\text{ cm}}$, environ 18,8 cm.
\label{fig-cor-16-5}\figureGeo[corrige]{espace-05}
\end{corrige}\end{blocLivre}
\begin{blocLivre}{cor-16-6}\begin{corrige}{6}{Sphère, boule et grands cercles}\demarche{Comparer chaque distance au centre avec le rayon.}
A, B et C sont à 4 cm de O : ils appartiennent à la $\res{\text{sphère}}$, donc aussi à la boule. M est à 2 cm de O : il est $\res{\text{dans la boule, pas sur la sphère}}$. N est à 6 cm : $\res{\text{à l’extérieur}}$. [AB] est un diamètre : $\res{AB=8\text{ cm}}$. Le grand cercle passant par A et B a pour longueur $\res{8\pi\text{ cm}}$.
\label{fig-cor-16-6}\figureGeo[corrige]{espace-sphere-vocab}
\end{corrige}\end{blocLivre}
\begin{blocLivre}{cor-16-7}\begin{corrige}{7}{Une pyramide dans un cube}\demarche{Repérer la base, la hauteur [EA] et les angles droits des faces.}
La base ABCD est un carré. EAB et EAD sont des triangles $\res{\text{rectangles isocèles en A}}$, moitiés de faces du cube. La droite (BC) est perpendiculaire à la face ABFE, donc à (BE) : EBC est $\res{\text{rectangle en B}}$ ; de même EDC est rectangle en D. Le volume vaut $36\times6/3=\res{72\text{ cm}^3}$, soit $\res{1/3}$ du cube de $216\text{ cm}^3$.
\label{fig-cor-16-7}\figureGeo[corrige]{espace-pyramide-cube}
\end{corrige}\end{blocLivre}
\begin{blocLivre}{cor-16-8}\begin{corrige}{8}{Une section de pavé}\demarche{Reconnaître une section parallèle à une face.}
La section est un $\res{\text{rectangle de 8 cm sur 5 cm}}$, comme la base. Son aire vaut $\res{40\text{ cm}^2}$. La hauteur de coupe ne modifie pas cette section parallèle dans un pavé.
\label{fig-cor-16-8}\figureGeo[corrige]{espace-section-pave}
\end{corrige}\end{blocLivre}
\begin{blocLivre}{cor-16-9}\begin{corrige}{9}{Sections d’un cylindre}\demarche{Identifier la direction du plan par rapport à l’axe.}
Perpendiculairement à l’axe, la section est un $\res{\text{disque de rayon 3 cm}}$, d’aire $\res{9\pi\text{ cm}^2}$. Un plan contenant l’axe donne un $\res{\text{rectangle de 6 cm sur 10 cm}}$ : sa largeur est un diamètre de la base. Son aire vaut $\res{60\text{ cm}^2}$.
\label{fig-cor-16-9}\figureGeo[corrige]{espace-section-cylindre}
\end{corrige}\end{blocLivre}
\begin{blocLivre}{cor-16-10}\begin{corrige}{10}{Une section de sphère}\demarche{Passer par le triangle rectangle de la section.}
Pour M sur le cercle, OHM est rectangle en H. $HM^2=OM^2-OH^2=25-9=16$, donc $\res{HM=4\text{ cm}}$. Le disque a pour aire $\res{16\pi\text{ cm}^2}$.
\label{fig-cor-16-10}\figureGeo[corrige]{espace-section-sphere}
\end{corrige}\end{blocLivre}
\begin{blocLivre}{cor-16-11}\begin{corrige}{11}{Un réservoir assemblé}\demarche{Additionner les volumes puis convertir.}
Le cylindre vaut $2\pi$ m$^3$, la demi-boule $\frac12\times\frac43\pi=\frac{2\pi}3$ m$^3$. Le total est $8\pi/3$ m$^3$, soit $8000\pi/3$ L, environ $\res{8378\text{ L}}$. Les deux parties ne se recouvrent qu’à leur frontière.
\label{fig-cor-16-11}\figureGeo[corrige]{espace-reservoir}
\end{corrige}\end{blocLivre}
\begin{blocLivre}{cor-16-12}\begin{corrige}{12}{Une boule creusée}\demarche{Soustraire les volumes du plein et du creux.}
On soustrait les volumes : $V=\frac43\pi(5^3-2^3)=\frac43\pi\times117=\res{156\pi\text{ cm}^3}$. Il faut soustraire des cubes, pas calculer le volume d’une boule de rayon $5-2$.
\label{fig-cor-16-12}\figureGeo[corrige]{espace-cavite}
\end{corrige}\end{blocLivre}
\begin{blocLivre}{cor-16-13}\begin{corrige}{13}{Un cône coupé}\demarche{Utiliser les dimensions proportionnelles de la section, puis soustraire.}
La section est une réduction de la base de coefficient $6/12=1/2$ : un $\res{\text{disque de rayon 3 cm}}$. Le petit cône a pour rayon 3 cm et hauteur 6 cm : $V=\pi\times9\times6/3=\res{18\pi\text{ cm}^3}$. Le grand cône vaut $\pi\times36\times12/3=144\pi$ cm$^3$ ; le tronc vaut $\res{126\pi\text{ cm}^3}$.
\label{fig-cor-16-13}\figureGeo[corrige]{espace-cone-coupe}
\end{corrige}\end{blocLivre}
\begin{blocLivre}{cor-16-14}\begin{corrige}{14}{Une section de pyramide}\demarche{Calculer le coefficient de réduction à partir des hauteurs.}
La section est une réduction de la base de coefficient $3/9=1/3$ : un $\res{\text{carré de côté 2 cm}}$, d’aire $\res{4\text{ cm}^2}$. La petite pyramide a pour hauteur 3 cm : $V=4\times3/3=\res{4\text{ cm}^3}$.
\label{fig-cor-16-14}\figureGeo[corrige]{espace-section-pyramide}
\end{corrige}\end{blocLivre}
\begin{blocLivre}{cor-16-15}\begin{corrige}{15}{Le cylindre d’Archimède}\demarche{Écrire les deux volumes avec le même rayon, puis simplifier.}
Le cylindre vaut $\pi\times9\times6=54\pi$ cm$^3$ et la boule $36\pi$ cm$^3$. Le rapport vaut $36/54=\res{2/3}$. Pour un rayon $r$ : le cylindre vaut $\pi r^2\times2r=2\pi r^3$ et la boule $\frac43\pi r^3$, d’où $\res{\frac{4/3}{2}=\frac23}$ pour tout $r$. Archimède avait découvert ce résultat.
\label{fig-cor-16-15}\figureGeo[corrige]{espace-archimede}
\end{corrige}\end{blocLivre}
